\documentclass[12pt]{article}
\usepackage{fullpage}
\usepackage{amsmath,amssymb,amsthm}

\newtheorem{lemma}{Lemma}
\newtheorem{proposition}{Proposition}
\newcommand{\R}{\mathbb R}
\newcommand{\E}{\mathbb E}
\newcommand{\1}{\mathbf 1}
\newcommand{\cE}{\mathcal E}
\newcommand{\cB}{\mathcal B}
\newcommand{\Var}{\operatorname{Var}}
\newcommand{\Lip}{\operatorname{Lip}}

\begin{document}

\begin{center}
{\bf Uniqueness of the invariant measure for the skew heat equation}
\end{center}

\section*{Problem statement and interpretation}

I use the standard interpretation of the distributional drift in the
statement: the solution is the Markov solution obtained as the limit of the
regularisations \(G_{1/n}\), equivalently the unique ABLM martingale/mild
solution.  With \(H_+(r):=\mathbf 1_{\{r>0\}}\), the drift
\(\delta_0\) is
\[
        -\frac12 \frac{d}{dr}(-2H_+(r)).
\]
Thus the equation is the singular gradient stochastic heat equation with
bounded potential \(-2H_+\).  The Dirichlet-form construction of the
corresponding skew stochastic heat equation (Bounebache--Zambotti) gives a
Markov solution with the same regularising approximations; by the assumed
weak uniqueness it is the same semigroup as \(P_t\).  The proof below uses
only the consequences of this identification that are stated explicitly in
Lemma 1.

Let \(E=C_0([0,1])\), \(H=L^2(0,1)\), and let \(\gamma\) be the law on
\(E\) of the standard Brownian bridge from \(0\) to \(0\).  Equivalently,
\(\gamma\) is the invariant law of the linear Dirichlet stochastic heat
equation.  Define
\[
        \Lambda(v):=\int_0^1 H_+(v(x))\,dx,\qquad
        \pi(dv):=Z^{-1}e^{2\Lambda(v)}\,\gamma(dv).
\]
Since \(0\leq\Lambda\leq1\), the density \(d\pi/d\gamma\) is bounded above
and below by positive constants.  In particular \(\pi\) is a probability
measure on \(E\) and \(\pi\sim\gamma\).  The value chosen for \(H_+(0)\) is
irrelevant for \(\gamma\) (and hence for \(\pi\)), because
\[
  \E_\gamma\!\int_0^1\mathbf 1_{\{\beta(x)=0\}}\,dx
  =\int_0^1\mathbb P(\beta(x)=0)\,dx=0.
\]

\begin{lemma}[reversibility and absolute continuity]\label{lem:DF}
For the semigroup \(P_t\) of the SPDE in the statement, the measure
\(\pi\) is invariant and symmetrizing:
\[
       \int_E g\,P_t f\,d\pi=\int_E f\,P_t g\,d\pi
       \qquad (f,g\in L^2(\pi),\ t\ge0).
\]
Moreover, for every \(\lambda>0\) the resolvent
\[
       R_\lambda(x,A):=\int_0^\infty e^{-\lambda t}P_t(x,A)\,dt
\]
satisfies \(R_\lambda(x,\cdot)\ll \pi\) for every \(x\in E\).
\end{lemma}

\begin{proof}
Put \(B_n(r):=\int_{-\infty}^r G_{1/n}(s)\,ds\).  For the regularised
equation with drift \(G_{1/n}=B_n'\), the formal generator on smooth
cylinder functions is
\[
 L_nF(v)=\frac12\operatorname{Tr}D^2F(v)
       +\langle {\tfrac12}\Delta_D v+G_{1/n}(v),DF(v)\rangle_H .
\]
The Gaussian integration-by-parts formula for the Brownian bridge gives
\[
   \int G\,L_nF\,d\pi_n
   =-\frac12\int\langle DF,DG\rangle_H\,d\pi_n
   =\int F\,L_nG\,d\pi_n,
\]
where
\[
   \pi_n(dv)=Z_n^{-1}\exp\!\left(2\int_0^1 B_n(v(x))\,dx\right)\gamma(dv).
\]
Hence \(\pi_n\) is the reversible invariant law of the regularised
gradient equation.  Since \(B_n(r)\to H_+(r)\) for \(r\neq0\) and
\(0\le B_n\le1\), dominated convergence gives \(\pi_n\Rightarrow\pi\).

The closed limiting form is
\[
       \cE(F,G)=\frac12\int_E\langle DF,DG\rangle_H\,d\pi,
       \qquad F,G\ \hbox{cylindrical},
\]
closed in \(L^2(\pi)\).  We now invoke the Dirichlet-form theorem for the
skew stochastic heat equation with bounded variation potential
\(f\) (Bounebache--Zambotti): for
\[
       \nu_f(dv)=Z_f^{-1}\exp\!\left(-\int_0^1 f(v(x))\,dx\right)\gamma(dv)
\]
the form \(\frac12\int\langle DF,DG\rangle\,d\nu_f\) is quasi-regular,
its associated process solves the corresponding local-time/skew heat
equation, and its resolvent has kernels absolutely continuous with respect
to \(\nu_f\).  Taking \(f=-2H_+\) gives exactly the above \(\pi\).  The
Fukushima decomposition of this process is the weak form of
\[
        \partial_t u=\tfrac12\partial_{xx}u+\delta_0(u)+\dot W
\]
with Dirichlet boundary conditions; for smooth \(f_n=-2B_n\) this is just
the regularised equation with drift \(G_{1/n}\), and the singular term is
the limit \(B_n'\to\delta_0\).  By the assumed uniqueness in law of the
ABLM solution on \(E\), this Dirichlet-form Markov process has semigroup
\(P_t\).  Thus \(P_t\) inherits symmetry, invariance of \(\pi\), and the
resolvent absolute-continuity property.
\end{proof}

\begin{lemma}[ergodic abstract criterion]\label{lem:abstract}
Let \(P_t\) be a Markov semigroup on a measurable space and let \(\pi\) be
an invariant probability measure.  Assume:
\begin{enumerate}
\item \(P_t\) is symmetric on \(L^2(\pi)\);
\item \(P_t\) has an \(L^2(\pi)\) spectral gap, i.e. for some \(c>0\),
\[
       \|P_t f-\pi(f)\|_{L^2(\pi)}
       \le e^{-ct}\|f-\pi(f)\|_{L^2(\pi)},\qquad f\in L^2(\pi);
\]
\item for some \(\lambda>0\), \(R_\lambda(x,\cdot)\ll\pi\) for every
\(x\).
\end{enumerate}
Then \(P_t\) has at most one invariant probability measure, namely \(\pi\).
\end{lemma}

\begin{proof}
Let \(\eta\) be invariant.  If \(\pi(A)=0\), then by invariance
\[
  \eta(A)=\lambda\int R_\lambda(x,A)\,\eta(dx)=0,
\]
so \(\eta\ll\pi\).  Write \(h=d\eta/d\pi\).  Symmetry implies that the
\(L^1(\pi)\)-adjoint of \(P_t\) is again \(P_t\); hence invariance of
\(\eta\) gives \(P_t h=h\) in \(L^1(\pi)\).

Fix \(a>0\) and set \(h_a=h\wedge a\).  Since \(r\mapsto r\wedge a\) is
concave and \(P_t\) is Markov,
\[
      h_a=(P_t h)\wedge a\ge P_t(h\wedge a)=P_t h_a .
\]
The two sides have the same \(\pi\)-integral because \(\pi\) is invariant;
therefore \(P_t h_a=h_a\) \(\pi\)-a.s.  Now \(h_a\in L^2(\pi)\), and the
spectral gap yields
\[
   \Var_\pi(h_a)=\Var_\pi(P_t h_a)\le e^{-2ct}\Var_\pi(h_a).
\]
For any \(t>0\) this forces \(\Var_\pi(h_a)=0\), so \(h\wedge a\) is
constant for every \(a>0\).  Hence \(h\) itself is constant \(\pi\)-a.s.;
because \(\int h\,d\pi=1\), \(h=1\).  Thus \(\eta=\pi\).
\end{proof}

It remains only to verify the spectral gap required in Lemma
\ref{lem:abstract}.  The Brownian bridge Gaussian measure satisfies the
Gaussian Poincare inequality
\[
   \Var_\gamma(F)\le C_0\int_E \|DF\|_H^2\,d\gamma
\]
on smooth cylinder functions (and hence on the closed domain).  Since
\(d\pi/d\gamma\) is bounded above and below, there is \(C<\infty\) such
that
\[
        \Var_\pi(F)\le C\,\cE(F,F),
        \qquad F\in D(\cE).
\]
By the spectral theorem for the symmetric semigroup associated with
\(\cE\), this Poincare inequality is exactly an \(L^2(\pi)\) spectral gap.

Combining this spectral gap with Lemma \ref{lem:DF}, all hypotheses of
Lemma \ref{lem:abstract} hold.  Therefore every invariant probability
measure of \(P_t\) equals
\[
   \boxed{\;
   \pi(dv)=Z^{-1}\exp\!\left(2\int_0^1 \mathbf 1_{\{v(x)>0\}}\,dx\right)
           \gamma(dv)\; } .
\]
In particular, \(P_t\) has at most one invariant probability measure.

\begin{thebibliography}{9}
\bibitem{ABLM}
S. Athreya, O. Butkovsky, K. L\^e and L. Mytnik,
\emph{Well-posedness of stochastic heat equation with distributional drift
and skew stochastic heat equation}, Comm. Pure Appl. Math. 77 (2024),
2708--2777.

\bibitem{BZ}
S. K. Bounebache and L. Zambotti,
\emph{A skew stochastic heat equation}, J. Theoret. Probab. 27 (2014),
168--201.
\end{thebibliography}

\end{document}
